% !TeX root = ../tesi.tex
% !TeX encoding = UTF-8
% !TeX spellcheck = it_IT

\chapter{Caratteristiche visive dell'immagine intera}
\label{chap:caratteristiche-visive}

In questo capitolo descrivo le caratteristiche visive che affianco al punteggio TruFor e al testo delle regioni sospette. Chiarisco il ruolo di questo ramo nella pipeline (Sezione~\ref{sec:feature-ruolo}), definisco le 19 misure (Sezione~\ref{sec:materiali-feature-visive}) e descrivo come le estraggo e le associo alle osservazioni (Sezione~\ref{subsec:metodi-feature-estrazione}). Concludo con i limiti della rappresentazione (Sezione~\ref{subsec:metodi-feature-limiti}).

\section{Ruolo delle caratteristiche visive nella pipeline}
\label{sec:feature-ruolo}

Affianco alla descrizione delle regioni sospette una rappresentazione dell'aspetto complessivo dell'immagine. Misuro proprietà come luminosità, contrasto, nitidezza, colore e composizione, introdotte nella Sezione~\ref{sec:background-feature-visive}. Uso queste informazioni per studiare se le proprietà visive sono associate all'engagement e al punteggio globale di TruFor.

Questo ramo riceve l'immagine originale selezionata dal filtraggio e comprende fotografia, scritte e componenti grafiche. Le descrizioni del Capitolo~\ref{chap:localizzazione-descrizione} dipendono invece dal contenuto conservato dalla maschera. Il calcolo delle caratteristiche visive usa direttamente i pixel dell'immagine intera e non dipende dalle regioni selezionate da TruFor. Ottengo così due rappresentazioni da confrontare nelle analisi successive: il contenuto descritto nelle regioni sospette e le proprietà visive complessive.

\section{Definizione delle caratteristiche visive}
\label{sec:materiali-feature-visive}

Rappresento ogni immagine $i$ con un vettore $\mathbf{f}_i=(f_{i1},\ldots,f_{i19})$ di 19 caratteristiche. Calcolo queste misure sull'immagine intera: comprendono quindi fotografia, testo e componenti grafiche. Le raggruppo nelle famiglie della Tabella~\ref{tab:feature-famiglie}.

\begin{table}[htbp]
    \centering
    \small
    \caption{Le cinque famiglie di caratteristiche estratte da ciascuna immagine.}
    \label{tab:feature-famiglie}
    \begin{tabular}{@{}>{\raggedright\arraybackslash}p{0.25\linewidth}r>{\raggedright\arraybackslash}p{0.52\linewidth}@{}}
        \toprule
        Famiglia & Numerosità & Proprietà \\
        \midrule
        Tono ed esposizione & 3 & Luminosità, contrasto e gamma dinamica. \\
        Fuoco e nitidezza & 4 & Nitidezza, limpidezza, sfocatura e profondità di campo. \\
        Colore & 4 & Vivacità, saturazione, calore e tonalità. \\
        Distribuzione delle intensità & 3 & Curtosi, asimmetria ed entropia. \\
        Composizione e struttura & 5 & Rapporto tra larghezza e altezza, simmetria, dominanza diagonale, regola dei terzi e densità dei bordi. \\
        \bottomrule
    \end{tabular}
\end{table}

\newpage
Per le definizioni considero una singola immagine di larghezza $W$ e altezza $H$. Indico con $R$, $G$ e $B$ i canali rosso, verde e blu, con intensità tra 0 e 255. Uso $\mu(A)$ e $\sigma(A)$ per la media e la deviazione standard degli elementi di una matrice $A$, calcolate sui suoi pixel.

Per elaborare le immagini ricorro a \emph{OpenCV} (Open Source Computer Vision Library), una libreria con funzioni per leggere, trasformare e analizzare immagini~\cite{opencv_library}.

Per descrivere la luminosità e le variazioni di intensità ricavo dai tre canali di colore una sola intensità per pixel. Uso due matrici: $L$ contiene una media pesata dei canali RGB; $Q$ contiene l'immagine in scala di grigi prodotta da OpenCV, con intensità intere da 0 a 255.

Assegno pesi diversi ai tre canali per approssimare il loro diverso contributo alla luminosità percepita: il verde contribuisce maggiormente, seguito dal rosso e dal blu. Calcolo quindi $L$:
\begin{equation}
    L=0{,}299R+0{,}587G+0{,}114B.
    \label{eq:feature-intensita}
\end{equation}
Uso $L$ per luminosità, contrasto e gamma dinamica; uso $Q$ per le misure di dettaglio, distribuzione delle intensità e struttura. Interpreto $L$ come un'approssimazione della luminosità~\cite{itu_bt601}.

Nelle sezioni seguenti impiego queste notazioni per definire le singole caratteristiche visive. Riporto tra parentesi i nomi delle colonne con cui salvo le caratteristiche nel \path{dataset_filtered.csv}.

\subsection{Tono ed esposizione}
\label{subsec:feature-tono}

La luminosità (\texttt{brightness}) è $\mu(L)/255$ e il contrasto (\texttt{contrast}) è $\sigma(L)/255$.

Indico con $P_q(L)$ il percentile di ordine $q$ delle intensità di $L$, cioè la soglia sotto cui si trova il $q\%$ dei pixel. Calcolo la gamma dinamica robusta (\texttt{dynamic\_range}) come segue:
\begin{equation}
    \texttt{dynamic\_range}=\frac{P_{99}(L)-P_1(L)}{255}.
    \label{eq:feature-gamma-dinamica}
\end{equation}
Escludendo le code dell'istogramma riduco il peso delle intensità estreme isolate. Questa misura indica quanto variano i livelli di luminosità nell’immagine, non quanta differenza tra luci e ombre riesce a registrare la fotocamera~\cite{1640788}.

\subsection{Fuoco e nitidezza}
\label{subsec:feature-fuoco}

Uso il \emph{Laplaciano}, un operatore basato sulle derivate seconde dell'intensità, per evidenziare le variazioni locali. Indico con $\Delta Q$ la sua risposta discreta.

La nitidezza (\texttt{sharpness}) è la varianza $\sigma^2(\Delta Q)$. Valori elevati indicano una maggiore variabilità della risposta, che può dipendere da dettagli, contorni o rumore~\cite{903548}.

Per misurare la limpidezza (\texttt{clearness}) confronto l'immagine in scala di grigi $Q$ con una sua versione sfocata. Ottengo quest'ultima applicando un filtro gaussiano, che attenua le variazioni tra pixel vicini. Indico con $G_2(Q)$ l'immagine filtrata con deviazione standard di 2 pixel e calcolo:
\begin{equation}
    \texttt{clearness}=\frac{\sqrt{\mu\bigl((Q-G_2(Q))^2\bigr)}}{255}.
\end{equation}
La misura cresce quanto più l'immagine differisce dalla sua versione sfocata: indica quindi quanto sono marcate le variazioni di intensità rimosse dal filtro~\cite{PERTUZ20131415}.

Per misurare la sfocatura (\texttt{blur}) sfrutto il principio che un'immagine già sfocata perde poco dettaglio quando la sfoco ulteriormente. Sfoco $Q$ separatamente in orizzontale e in verticale, usando filtri che calcolano la media su 9 pixel. Per ciascuna direzione misuro quanto si riducono le differenze di intensità tra pixel vicini rispetto a quelle iniziali. Una riduzione piccola produce un indice di sfocatura alto. Conservo il maggiore dei due indici, compreso tra 0 e 1 e crescente con il grado di sfocatura~\cite{creteroffet:hal-00232709}.

La profondità di campo (\texttt{depth\_of\_field}) è un indicatore del rapporto di dettaglio tra centro e periferia di una foto. Definisco il centro $C$ come il rettangolo che occupa il terzo centrale di ciascuna dimensione e la periferia $P$ come tutti i pixel restanti. Calcolo
\begin{equation}
    \texttt{depth\_of\_field}=
    \frac{\sigma^2((\Delta Q)_C)}{\sigma^2((\Delta Q)_P)+10^{-6}}.
\end{equation}
Il termine $10^{-6}$ rende il rapporto definito anche con una periferia uniforme. Un rapporto elevato indica più dettaglio al centro; l'interpretazione come soggetto a fuoco su sfondo sfocato richiede che il soggetto occupi quella regione~\cite{10.1007/11744078_23}.

\subsection{Colore}
\label{subsec:feature-colore}

Per la vivacità (\texttt{colorfulness}) uso la misura di Hasler e Süsstrunk. Definisco i canali di opposizione $U=R-G$ e $V=(R+G)/2-B$ e calcolo:
\begin{equation}
    \texttt{colorfulness}=
    \sqrt{\sigma^2(U)+\sigma^2(V)}+
    0{,}3\sqrt{\mu(U)^2+\mu(V)^2}.
\end{equation}
La misura vale zero quando i tre canali coincidono. Mantengo la scala originale, senza normalizzarla~\cite{10.1117/12.477378}.

Converto l'immagine nel formato HSV (Hue, Saturation, Value). Nella rappresentazione OpenCV a 8 bit la tonalità varia da 0 a 179 e la saturazione e il valore da 0 a 255~\cite{opencv_library}.

La saturazione (\texttt{saturation}) è la media del relativo canale divisa per 255.

La tonalità media (\texttt{hue\_mean}) è la media aritmetica del canale di tonalità divisa per 179.

Il calore (\texttt{warmth}) è $\mu\bigl((R-B)/(R+B+10^{-6})\bigr)$. Il segno positivo indica una prevalenza del rosso sul blu, quello negativo la prevalenza opposta. Questa definizione produce un indice cromatico compreso tra $-1$ e $1$~\cite{Jonauskaite2024}.

\subsection{Distribuzione delle intensità}
\label{subsec:feature-distribuzione}

Descrivo la forma dell'istogramma di $Q$ attraverso i momenti centrali $m_k=\mu\bigl((Q-\mu(Q))^k\bigr)$. L'asimmetria (\texttt{skewness}) è $m_3/m_2^{3/2}$. La curtosi (\texttt{kurtosis}) è $m_4/m_2^2-3$: uso quindi l'eccesso di curtosi, che per una distribuzione normale vale zero. Calcolo entrambi i rapporti con le impostazioni predefinite di SciPy, senza correzione del bias~\cite{2020SciPy-NMeth}.

L'\emph{entropia di Shannon} (\texttt{shannon\_entropy}) riassume la dispersione delle frequenze dei livelli di grigio. Se $p_v$ è la frazione di pixel con intensità $v$, definisco
\begin{equation}
    \texttt{shannon\_entropy}=-\sum_{v:p_v>0}p_v\log_2 p_v.
\end{equation}
L'entropia è espressa in bit, tra 0 e 8 per 256 livelli. Vale zero per un'immagine uniforme e cresce quando le frequenze si distribuiscono su più livelli. L'istogramma ignora la posizione dei pixel: immagini con disposizioni spaziali diverse possono avere la stessa entropia~\cite{6773024}.

\subsection{Composizione e struttura}
\label{subsec:feature-composizione}

L'\texttt{aspect\_ratio} è $W/H$: valori superiori a 1 indicano un formato orizzontale, inferiori a 1 un formato verticale.

La simmetria (\texttt{symmetry}) confronta la metà sinistra $Q_S$ con la metà destra riflessa orizzontalmente $Q_D^{\leftrightarrow}$. Calcolo $1-\mu(|Q_S-Q_D^{\leftrightarrow}|)/255$, escludendo la colonna centrale quando la larghezza è dispari. L'indice appartiene a $[0,1]$ ed è massimo quando le due metà coincidono nelle intensità.

Per le due misure successive calcolo gli operatori di Sobel di dimensione $3\times3$ per misurare le variazioni di intensità in orizzontale ($g_x$) e in verticale ($g_y$). Le combino nel modulo del gradiente $M=\sqrt{g_x^2+g_y^2}$, che indica quanto è forte la variazione in ciascun pixel~\cite{996}.

La dominanza diagonale (\texttt{diagonal\_dominance}) misura quanto prevalgono le variazioni orientate lungo le diagonali. Considero le direzioni vicine a $45^\circ$ e $135^\circ$, con una tolleranza di $22{,}5^\circ$. Sommo le intensità del gradiente $M$ nei pixel con queste direzioni e divido per la somma di $M$ su tutta l'immagine. Ottengo una quota compresa tra 0 e 1, che cresce con la presenza di variazioni diagonali, pesate per la loro forza.

La regola dei terzi (\texttt{rule\_of\_thirds}) misura la concentrazione del gradiente vicino alle quattro intersezioni di una griglia $3\times3$. Costruisco attorno ad ogni intersezione un disco di raggio pari al $10\%$ della diagonale dell'immagine. Indico con $D$ l'unione dei quattro dischi e con $M_D$ le intensità del gradiente nei pixel di $D$ e calcolo:
\begin{equation}
    \texttt{rule\_of\_thirds}=\frac{\mu(M_D)}{\mu(M)}.
    \label{eq:feature-regola-terzi}
\end{equation}
Un rapporto superiore a 1 indica una maggiore concentrazione di variazioni in quelle zone. Se il gradiente è nullo su tutta l'immagine, pongo la caratteristica a zero.

La complessità visiva (\texttt{visual\_complexity}) è la frazione di pixel rilevati come bordi dall'algoritmo di Canny, applicato a $Q$ con soglie 100 e 200. La densità dei bordi è compresa tra 0 e 1~\cite{4767851}.

\section{Estrazione e aggiornamento della tabella}
\label{subsec:metodi-feature-estrazione}

Dopo l'inferenza di TruFor eseguo \path{scripts/features.py}, che associa a ciascuna osservazione il vettore definito nella Sezione~\ref{sec:materiali-feature-visive}.

Leggo le immagini selezionate dalla cartella \path{scripts/temp/images/} attraverso OpenCV~\cite{opencv_library} e calcolo le caratteristiche visive sull'intera immagine con OpenCV, NumPy~\cite{harris2020array}, SciPy~\cite{2020SciPy-NMeth} e scikit-image~\cite{scikit-learn}. Applico le normalizzazioni previste dalle singole definizioni. 

Nella configurazione corrente imposto \texttt{FEATURE\_RESIZE\_LONG\_SIDE} a 0, mantengo quindi la risoluzione originale dei file letti. Lo script consente di fissare un limite al lato maggiore: in quel caso riduce le immagini che lo superano mantenendo le proporzioni, con interpolazione per area. Il rapporto d'aspetto viene calcolato prima dell'eventuale ridimensionamento. Questa impostazione è indipendente dal limite di risoluzione adottato per TruFor.

Salvo le caratteristiche nelle scale definite nella Sezione~\ref{sec:materiali-feature-visive}. Associo il vettore al campo \texttt{filename}, ottenuto dal nome del file senza estensione, e aggiorno \path{dataset_filtered.csv}. Se le colonne sono già presenti, sostituisco i valori precedenti. Elimino le righe con valori mancanti.

Gestisco separatamente i casi degeneri. Salto i file che non riesco a leggere. Sostituisco con zero i valori non finiti, come asimmetria e curtosi di un'immagine uniforme. Per una direzione priva di variazioni assegno zero al relativo indice di sfocatura; assegno zero anche alla dominanza diagonale e alla regola dei terzi quando il gradiente è nullo. Questi zeri sono convenzioni numeriche: in particolare, l'assenza di variazioni rende la sfocatura non valutabile attraverso il confronto adottato.

\section{Limiti della rappresentazione}
\label{subsec:metodi-feature-limiti}

L'interpretazione dipende dal contenuto e dalla scala dei file. Le finestre di filtraggio sono espresse in pixel: risoluzione, compressione e rumore possono modificare nitidezza, sfocatura e bordi. Testo al centro o sfondi uniformi possono aumentare il rapporto centro-periferia indipendentemente dalla messa a fuoco.

La tonalità ha natura circolare: valori vicini ai due estremi possono rappresentare colori simili. La media aritmetica trascura questa proprietà e include pixel poco saturi. Entropia e densità dei bordi descrivono frequenze e transizioni di intensità, da distinguere dalla complessità semantica della scena.

Le misure aggregano regioni autentiche ed eventualmente alterate. Le associazioni con il punteggio TruFor $t_i$ descrivono il comportamento del rilevatore. Per verificare se queste caratteristiche siano associate anche a manipolazioni effettive servono annotazioni indipendenti. Le correlazioni con il punteggio di engagement $c_i$ mostrano associazioni, che possono dipendere anche dall’account, dall’argomento e dalle condizioni di pubblicazione.

